\documentclass[11pt,a4paper]{article}
\usepackage[utf8]{inputenc}
\usepackage[T1]{fontenc}
\usepackage[margin=0.75in]{geometry}
\usepackage{enumitem}
\usepackage{titlesec}
\usepackage{xcolor}
\usepackage{hyperref}

\definecolor{primarycolor}{RGB}{33,37,41}
\definecolor{secondarycolor}{RGB}{0,102,204}
\definecolor{accentcolor}{RGB}{108,117,125}
\hypersetup{colorlinks=true,linkcolor=secondarycolor,urlcolor=secondarycolor}
\titleformat{\section}{\large\bfseries\color{primarycolor}}{}{0em}{}
\titlespacing*{\section}{0pt}{1.1em}{0.4em}
\setlist[itemize]{leftmargin=1.3em,itemsep=0.2em,topsep=0.2em}
\pagestyle{empty}
\setlength{\parindent}{0pt}
\setlength{\parskip}{0.25em}

\begin{document}

\begin{center}
{\LARGE\bfseries\color{primarycolor} MARTA KOWAL-BRANDT}\\
QA Automation Engineer transitioning to DevOps / Platform Engineering\\
Wroclaw, Poland | \href{mailto:marta.kowalbrandt@example.com}{marta.kowalbrandt@example.com} | +48 555 303 303 | \url{gitlab.com/mkb-example}
\end{center}

\section*{PROFESSIONAL SUMMARY}
QA automation engineer with 7 years of experience, including two years focused primarily on CI/CD pipelines, test environments and build tooling. Hands-on experience owning GitLab CI pipelines for three product teams and building Helm-based ephemeral environments on a shared Kubernetes cluster. Personal infrastructure work includes a k3s homelab, Terraform-provisioned Hetzner VPS infrastructure and introductory AWS learning. Seeking to move into a DevOps or platform engineering role.

\section*{TECHNICAL SKILLS}
CI/CD and containers: GitLab CI, Jenkins, Docker, Docker Compose\\
Infrastructure and platforms: Kubernetes (namespace-level access; k3s homelab), Helm, Terraform (personal projects), Ansible (basic), Linux\\
Monitoring and data: Grafana, Prometheus, Loki, PostgreSQL, SQL\\
Programming and testing: Python, Bash, Java (test code), pytest, Playwright, Selenium, Postman\\
Cloud learning: AWS (EC2 and S3 basics)\\
Other: Git, Argo CD

\section*{PROFESSIONAL EXPERIENCE}
\textbf{QA Automation Engineer (Test Infrastructure) | Odrava Health Software, Wroclaw | 2021-present}
\begin{itemize}[label=-]
\item Own GitLab CI pipelines for three product teams (~25 developers), covering builds, unit, integration and end-to-end test stages, and artifact publishing.
\item Reduced average pipeline time from 38 to 14 minutes by parallelising test jobs, caching dependencies and separating end-to-end tests into their own stage; monthly CI runner costs fell by about 30\%.
\item Built merge-request ephemeral environments with Docker Compose and later Helm on the company's shared Kubernetes cluster, including application services, PostgreSQL and seeded data; environments are removed after merge or 48 hours of inactivity.
\item Maintain a pytest and Playwright end-to-end suite of about 900 tests.
\item Built a flaky-test dashboard using Python, PostgreSQL, Grafana and the GitLab API; the flaky rate fell from 7\% to below 1\% after teams addressed the top offenders.
\item Wrote Bash and Python tooling to seed PostgreSQL test data and documented how the CI system works for team onboarding.
\item Mentored two QA interns in 2023.
\end{itemize}

\textbf{QA Engineer | Brightmill Software House, Wroclaw | 2019-2021}
\begin{itemize}[label=-]
\item Tested web and mobile applications through manual and automated testing for client projects.
\item Introduced a Java and Selenium automated regression suite for the largest client project.
\item Configured Jenkins jobs to run nightly regression tests.
\end{itemize}

\textbf{Customer Support Specialist (2nd line) | Brightmill Software House | 2018-2019}
\begin{itemize}[label=-]
\item Investigated customer issues using logs and databases and reproduced bugs for developers.
\end{itemize}

\section*{SELECTED INFRASTRUCTURE PROJECT}
\textbf{Homelab k3s Cluster with GitOps | Personal | 2024-2025}
\begin{itemize}[label=-]
\item Built a three-node k3s cluster on Raspberry Pi 5 devices, with applications deployed from a Git repository using Argo CD.
\item Set up Prometheus, Grafana and Loki for monitoring, plus Traefik ingress and cert-manager for TLS.
\item Provisioned infrastructure for a small Hetzner VPS with Terraform; used Ansible at a basic level.
\end{itemize}

\section*{EDUCATION}
BEng Computer Engineering | Silesian Lowlands University of Technology | 2018

\section*{CERTIFICATIONS AND PROFESSIONAL DEVELOPMENT}
ISTQB Foundation Level | 2019\\
Preparing for the Certified Kubernetes Administrator (CKA) exam, scheduled for January 2027

\section*{LANGUAGES}
Polish (native), English (C1), German (A2)

\end{document}